\section{E5: credit propagation}
\label{sec:credit}

\lead{The prediction.} Freezing is what makes the scheme greedy, and clause inclusion is a
credit path that does not need gradients. If layer-2 feedback can reach layer~1, the
saturation problem goes away by construction --- layer~1 is no longer optimising a proxy, it
is optimising the thing layer~2 wants.

This is the idea that changes the proposal most, and it is the one that did not survive first
contact with an implementation. Writing the rule out honestly produced a sequence of failures,
each of which had to be fixed before the next became visible: no forgetting term, then a
forgetting term that erases layer~1, then an event ratio nothing in the loop chooses, then no
events at all, and finally --- once there are events --- drift. The ladder
(Table~\ref{tab:credit}) is more informative than any single arm, and it is the substance of
this section.

\begin{table}[h]
\centering\small
\caption{The credit ladder. \code{rate} is the fraction of credit events actually applied,
\code{warm} the number of frozen-layer-1 epochs before the credit path is switched on,
\code{Ia}/\code{Ib} the feedback events actually delivered to layer~1 over the whole run, and
\code{size}/\code{fire} layer~1's median clause size and firing rate at the end. \code{best}
is the best test accuracy over epochs, as everywhere else in this report; \code{final} is the
accuracy at the last epoch, and the two differ here because several of these arms do not
converge --- they collapse, and the best epoch is an early one.}
\label{tab:credit}
\pgfplotstabletypeset[
  col sep=comma, string type,
  every head row/.style={before row=\ttoprule, after row=\ttmidrule},
  every last row/.style={after row=\ttbotrule},
  columns={label,rate,warmup,ev_i,ev_ib,l1_size,l1_fire,acc,final},
  columns/label/.style={column name=\thead{arm}, column type={l}},
  columns/rate/.style={column name=\thead{rate}, column type={r}},
  columns/warmup/.style={column name=\thead{warm}, column type={r}},
  columns/ev_i/.style={column name=\thead{Ia}, column type={r}},
  columns/ev_ib/.style={column name=\thead{Ib}, column type={r}},
  columns/l1_size/.style={column name=\thead{size}, column type={r}},
  columns/l1_fire/.style={column name=\thead{fire (\%)}, column type={r}},
  columns/acc/.style={column name=\thead{best (\%)}, column type={r}},
  columns/final/.style={column name=\thead{final (\%)}, column type={r}},
]{data/credit_rungs.csv}
\end{table}

\subsection{Rung 1: the rule as written has no forgetting term}

Taken literally --- ``a layer-2 clause that receives Type~I feedback credits Type~I feedback to
layer-1 clause $k$'' --- and applied to the case that has a position to point at, the rule
propagates Type~Ia and Type~II. Both of those \emph{add} literals. Type~Ia memorizes the
literals true in the drawn patch, and with \code{boost\_true\_positive} it does so with
probability~1, so a single event includes every literal true in that patch at once. Type~II
pushes false literals towards inclusion. Nothing in the rule ever removes one.

Layer~1 saturates inside the first epoch: clauses reach \lOneSizeMctmCreditIa\ literals,
\lOneDeadMctmCreditIa{}\% of channels go permanently dead and the median channel stops firing
altogether. The credit path then extinguishes itself, which is the part worth keeping. It
delivers \eventsMctmCreditIa\ Type~Ia events in total, all of them by epoch
\lastEventMctmCreditIa, and not one afterwards for the remaining \epochsRun-epoch run: a
layer-2 clause can only propagate through a positive channel literal, and there are almost no
positive channel activations left to include. The arm reaches \accMctmCreditIa{}\%, which is
what a mostly-dead layer~1 is worth.

\subsection{Rung 2: Type Ib has to be blamed, not sampled --- and it erases layer 1}

Type~I covers Ib as well as Ia, and Ib is the only term that removes literals. It also needs
no position, which makes it easy to propagate. The obvious implementation --- pick a random
included positive channel literal of the non-matching layer-2 clause and give layer-1 clause
$k$ a Type~Ib event --- is wrong, and instructively so: it blames layer-1 clauses that did
exactly what was asked of them. The correct rule asks which literal actually blocked the
match: take the clause at the patch it came closest to matching, and blame an included
positive channel literal that was \emph{false} there. That is the rule used here, and the
implementation is checked against it directly (\code{test\_ideas.py}).

It does not save layer~1; it destroys it the other way. Layer-2 clauses fail to match far more
often than they match, so blamed Type~Ib events outnumber Type~Ia by roughly three orders of
magnitude (\eventsIbMctmCreditIb\ against \eventsMctmCreditIb\ over the run), and layer~1 is
forgotten down to \lOneSizeMctmCreditIb\ literals per clause --- \lOneDeadMctmCreditIb{}\% of
channels dead, which is to say the representation is gone. The arm's best epoch is its first,
at \accMctmCreditIb{}\%, and it finishes at \finMctmCreditIb{}\%: exactly chance, because every
image now maps to the same all-zero feature map.

\subsection{Rung 3: balancing the two directions, and what that exposes}

Capping the Type~Ib budget at the Type~Ia rate actually delivered makes the equilibrium clause
size finite by construction, and the arm reaches \accMctmCreditBal{}\%. Layer~1 ends at
\lOneSizeMctmCreditBal\ literals and \lOneFireMctmCreditBal{}\% firing --- back where rung~1
was.

\begin{ttkey}
\lead{The credit path severs the loop that keeps a Tsetlin clause the right size.} In an
ordinary machine, a clause's ratio of Type~Ia to Type~Ib events \emph{is} its own match rate:
grow too specific, stop matching, receive Type~Ib, shrink. It is self-stabilising, and it is
why a trained layer settles at 77--85 literals rather than at 0 or 496. Credit propagation
replaces that ratio with \emph{layer 2's} match rate, which carries no information about
whether layer-1 clause $k$ is the right size. The three rungs are the three things that can
then happen: flood (rung~1, saturate at
\lOneSizeMctmCreditIa), drought (rung~2, erased to \lOneSizeMctmCreditIb), and an arbitrary
externally-chosen ratio (rung~3). There is no stable interior, because the quantity that
used to choose it is no longer in the loop.
\end{ttkey}

The second thing rung~3 exposes is bandwidth. The credit path's throughput is the rate at
which a layer-2 clause matches \emph{because of} a positive channel literal, and
Table~\ref{tab:credit} says what that is: \eventsMctmCreditBal\ Type~Ia events across
\epochsRun\ epochs and \nTrain\ images, against the $\nTrain{} \times \cOne$ per epoch a layer
trained directly receives. Layer~1 is not being trained through layer~2; it is being nudged a
few hundred times.

\subsection{Rung 4: warm start, and the bandwidth result at its sharpest}

Both problems have a candidate fix. An untrained layer~2 has no included positive channel
literals, so there is nothing to propagate through until it has some: a warm start --- ten
epochs of ordinary frozen training, then the credit path --- is the natural analogue of how
layer-wise pretraining is actually used. And the stabilising loop the credit path severed can
be replaced from outside, by running \ideaB's controller inside the training loop, which holds
layer~1 at a chosen firing rate no matter what the credit events do to it.

\begin{ttkey}
Warm-starting from the \emph{greedy} layer~1 delivers \eventsMctmCreditWarm\ credit events.
Not few --- none, across twenty epochs and \nTrain\ images. A layer firing on
\fireGreedyPre{}\% of positions never lets a layer-2 clause match on a positive channel
literal, so the rule has nothing to fire on and the arm is the frozen stack exactly
(\accMctmCreditWarm{}\% against \accMctm{}\%). Sweeping the credit step size on top of that
would be sweeping a multiplier on zero.
\end{ttkey}

Adding the controller is what gives the path traffic, and then the mechanism can finally be
judged on its own terms.

\begin{ttkey}
\code{mctm-credit-cal} holds layer~1 at \lOneFireMctmCreditCal{}\% firing and
\lOneSizeMctmCreditCal\ literals --- statistically the same layer that
\S\ref{sec:density}'s rate-only calibration produces and freezes
(\calibFireMctmCalibRTwoZero{}\%, \calibSizeMctmCalibRTwoZero\ literals), which reaches
\accMctmCalibRTwoZero{}\%. The credit path then delivers \eventsMctmCreditCal\ Type~Ia events,
and the arm reaches \accMctmCreditCal{}\%, finishing at \finMctmCreditCal{}\%. Turning the step
size up to $1.0$ gives \accMctmCreditCalPOne{}\% and down to $0.01$ gives
\accMctmCreditCalPZeroZeroOne{}\%.
\end{ttkey}

Layer~2 therefore \emph{ends} over the same representation as the frozen arm it should be
compared with, and the stack is nearly eight points worse. What it did not have is a
representation that stayed still while it learned: from epoch \creditWarmEpochs\ onwards every
layer-2 clause is a conjunction over channels whose meaning is being rewritten underneath it,
and the less that rewriting is allowed to happen, the better the stack does.
Figure~\ref{fig:credit} shows it over time --- the arms are identical through the warm start
and separate at the moment layer~1 starts to move.

\begin{figure}[h]
\centering
\begin{tikzpicture}
\begin{axis}[
  width=0.92\textwidth, height=6.0cm,
  xlabel={epoch}, ylabel={test accuracy (\%)},
  grid=major, grid style={ttline!40},
  legend style={font=\sffamily\scriptsize, draw=ttline, at={(0.5,-0.26)}, anchor=north,
                legend columns=4, column sep=1.1em},
  label style={font=\sffamily\small}, tick label style={font=\sffamily\footnotesize},
  every axis plot/.append style={very thick},
]
\addplot[ttslate] table[col sep=comma, x=epoch, y=test_acc] {data/credit_trace_mctm-credit-warm.csv};
\addplot[ttmoss] table[col sep=comma, x=epoch, y=test_acc] {data/credit_trace_mctm-credit-cal-p001.csv};
\addplot[ttamber] table[col sep=comma, x=epoch, y=test_acc] {data/credit_trace_mctm-credit-cal.csv};
\addplot[ttorange] table[col sep=comma, x=epoch, y=test_acc] {data/credit_trace_mctm-credit-cal-p1.csv};
\legend{no controller, step 0.01, step 0.1, step 1.0}
\end{axis}
\end{tikzpicture}
\caption{The warm-start arms. All four are ordinary frozen training for the first
\creditWarmEpochs\ epochs. The ``no controller'' arm never changes afterwards either, because
the credit path delivers no events at all over a layer~1 that fires on \fireGreedyPre{}\% of
positions.}
\label{fig:credit}
\end{figure}

\subsection{What the ladder actually shows}

Most of the rungs are implementation failures, and they are worth recording precisely because
the proposal's one-sentence statement of the rule hides every one of them: there is no
forgetting term in it, no account of which literal to blame, and no quantity that sets the
ratio between the two. An implementer meets all three before getting to the interesting
question.

The last two rungs are not implementation failures. \emph{No bandwidth} is a statement about
the setting --- the credit path can only carry traffic when layer-2 clauses match because of
positive channel literals, and a layer~1 firing on \fireGreedyPre{}\% of positions never lets
that happen. \emph{Drift} is a statement about the mechanism, and it is the one that does not
go away by fixing something: a frozen layer is what makes ``channel $k$'' mean something
stable enough for a layer-2 clause to be a rule about it, and credit propagation gives exactly
that up. \S\ref{sec:synthetic} runs the same code on a task where layer~2 clauses are small
and match often, and finds the same collapse.
